\section{Introduction}
\label{sec:intro}

Sleep staging is the foundation of every clinical sleep study. A technologist
scores each \SI{30}{\second} epoch of polysomnography into one of five stages
--- wake (W), N1, N2, N3 and REM --- and the resulting hypnogram drives the
diagnosis of insomnia, narcolepsy, sleep apnoea and parasomnias, as well as
the assessment of treatment response. Manual scoring is slow and only
moderately reliable, with inter-scorer agreement on N1 and on transition
epochs notoriously poor \citep{rosenberg2013}. Automatic staging has therefore
become a standard first pass, and recent single-channel models reach epoch
accuracy in the mid-80s on public cohorts
\citep{lawhern2018,phan2019seqsleepnet,perslev2021,phan2022sleepformer},
with rhythm, multi-scale and contrastive designs continuing to raise it
\citep{zhou2026bspc,fan2025bspc}.

The reviewing technologist is the part of that workflow a point prediction
does not serve. A stager that is right nine epochs in ten still returns a
single label for every epoch, including the ones it should not be trusted on,
so the reviewer either re-scores the whole night or trusts a number that was
never a statement about reliability. What the review step needs is a
per-epoch account of which stages remain plausible after the signal has been
read, so that attention can be spent where the model is genuinely undecided.
This is a signal-analysis question before it is a machine-learning one: the
decision has to be made from a short single-channel epoch, it has to be
consistent with the temporal structure of the hypnogram, and it has to be
usable on a patient the model has never seen.

Cohort averages hide two properties that matter at the point of review.

The first is \emph{subject shift}. A model is trained on a cohort and then
applied to a patient it has never seen. Sleep architecture varies enormously
between individuals --- total sleep time, N3 fraction, REM latency and the
amplitude and spectral content of the EEG all change with age, medication and
comorbidity --- and a model that is 85\% accurate on the cohort can be
markedly worse on a particular patient. A cohort mean does not tell a
clinician which patients those are. Per-epoch confidence scores (softmax
confidence, Monte-Carlo dropout, deep ensembles, evidential heads) produce a
number per epoch but no statement about how often the true stage is actually
inside whatever set of stages the model is willing to commit to.

The second is \emph{temporal dependence}. Sleep stages persist for minutes.
Scored epoch by epoch, a model's errors cluster, and any per-epoch
uncertainty estimate inherits that clustering: it is confident through a
stable stretch of N2 and uncertain precisely at the transitions that
clinicians care about. Per-epoch prediction sets can also flicker between
adjacent epochs even when the underlying signal has not changed, which is
clinically incoherent, because the hypnogram is a sequence rather than a bag
of epochs.

Conformal prediction (CP) \citep{vovk2005,angelopoulos2023gentle} is the
natural framework for the first problem. Given a calibration set, split CP
turns any model into a set-valued predictor with a finite-sample,
distribution-free guarantee that the true label lies in the set with
probability at least $1-\alpha$. In sleep staging this is exactly the
statement a clinician needs: ``with 90\% confidence, this epoch is one of
these two stages.'' CP has also been studied as a target for network
design and compression rather than only as a wrapper around a fixed model
\citep{zhao2025pruning}; we deliberately keep the backbone frozen, so that
every statement below is about the stager as it would be deployed and not
about a network retrained to make the sets smaller. Recent work has begun to
bring CP to EEG and to
sleep, and the closest results are sobering: split CP on EEG event
classification loses a large fraction of its coverage under patient shift,
and per-patient recalibration recovers most of it \citep{chatterjee2026}. Selective
prediction for sleep staging \citep{cover2026} takes a related but different
route --- it decides which epochs to accept, subject to an empirical risk
target, without producing a coverage-guaranteed set.

What is missing is (i) an evaluation of CP for sleep staging that reports the
distribution of coverage across patients rather than its mean, and (ii) a
procedure that delivers the per-patient statement a clinician needs while
respecting the information a deployment actually has: a frozen model, a
handful of scored nights from the cohort it was trained on, and labels that
arrive after each prediction because a technologist reviews the hypnogram
anyway.

This paper provides both, and is deliberately explicit about what the data
can and cannot support. We build on Sleep-EDF
\citep{kemp2000,goldberger2000} with subject-wise folds that partition the
cohort, so that every one of the 20 subjects is tested exactly once
and no test epoch shares a subject with any calibration or training epoch. We
evaluate two backbones, several nonconformity scores, and three families of
adaptation, and we separate the claims that carry a guarantee from those that
are merely measured.

\paragraph{The gap, and four findings}
\begin{enumerate}
\item \textbf{The marginal guarantee is not the clinical one.} At
      $\alpha=0.1$, pooled split CP reaches $0.898$ coverage, close to
      nominal, while $45\%$ of test subjects fall below the nominal level
      and the worst subject is covered at $0.774$. An epoch-level random
      split --- the protocol common in the machine-learning literature ---
      hides this entirely: calibration and test epochs then come from the
      same subjects, so a coverage figure computed per subject reflects the
      shared epochs rather than the subject and is not a deployment
      statistic.
\item \textbf{Score smoothing is what buys efficiency.} Running
      forward--backward recursions over the HMM implied by the training
      labels and scoring the \emph{posterior deficit}
      $\max_k \gamma_t(k)-\gamma_t(Y_t)$ cuts the set size by about forty per
      cent relative to APS on the raw softmax at comparable coverage
      ($1.53$ versus $2.54$ labels per epoch, at $0.889$ versus $0.898$), and raises temporal
      consistency from $0.68$ to $0.94$. Because a singleton set is an epoch
      the reviewer can accept, set size is also the review workload: the
      fraction of epochs whose set contains more than one stage falls from
      $0.76$ to $0.16$. The gain is not uniform across stages, and
      \Cref{sec:results-score} reports both the trade-off and the per-stage
      threshold that removes it at a modest cost in set size. The smoothing
      costs $O(LK^2)$ per recording and uses no labels at test time.
\item \textbf{One adaptation stream per patient shrinks the spread.} Running
      adaptive conformal inference (ACI) with an independent stream per
      patient, and choosing the step size by leave-one-subject-out simulation
      on the calibration subjects, raises worst-patient coverage from
      $0.774$ to $0.886$ without tuning on the test patients. The
      mechanism is visible in the per-patient coverage vectors: easy patients
      move down toward nominal, hard patients move up.
\item \textbf{Sequence-level set construction does not help here.} We derive
      two dynamic-programming constructors, one centred on the minimum-cost
      path and one projecting an exact block-level conformal set. We prove
      that the first reduces exactly to per-epoch CP on the deficit score when
      the transition weight is zero, and that its position-centred variant
      cannot return larger sets; at positive weight the measured sets do grow,
      and we report that as a measured effect rather than a proved one. We
      prove that the second is simultaneous by construction, and
      show empirically that it is vacuous on sleep blocks. We report these as
      negative results with proofs, because a reader who knows the
      literature will otherwise assume the constructions were not tried.
\end{enumerate}

\paragraph{Relation to prior work}
The procedure is assembled from parts that exist separately, and we say so
plainly rather than claim each part. \citet{gibbs2021} introduced ACI and
proved long-run coverage for a single stream under arbitrary shift; running
their recursion once per patient is what turns that pooled statement into a
per-patient one, and \Cref{cor:per-patient} records the inheritance under
their step-size conditions. \citet{angelopoulos2024} analyse decaying
schedules; the default schedule here fails their sufficient conditions, which
\Cref{sec:results-aci} states and quantifies rather than papers over.
\citet{barber2023} and \citet{tibshirani2019} extend CP beyond exchangeability,
but the resulting bounds are expressed in terms of the size of the shift or of
a known weighting function, neither of which is available at deployment, and
neither gives a per-patient statement. \citet{chatterjee2026} is the closest
empirical
predecessor: they measure the collapse of EEG coverage under patient shift and
repair it with per-patient recalibration, which needs labelled epochs from the
test patient before deployment. The warm-up variant studied in
\Cref{sec:method-sacp} is exactly their requirement, so we evaluate it rather
than assume it; the ACI variant replaces the labelled warm-up with labels that
arrive one at a time, which is what a review workflow already provides. On the
sleep side the models we use are standard and the contribution is not
architectural. What we contribute is the signal-processing layer that turns a
per-epoch posterior into a clinically readable object: a forward--backward
posterior computed from the epoch probabilities and the sleep-stage
transition structure, a nonconformity score that is monotone in the posterior
deficit and therefore inherits both, a calibration rule that respects the
patient as the unit of deployment, and an evaluation that reports the
quantities a review workflow actually consumes --- coverage per patient, set
size as review workload, and temporal consistency of the sets along the
hypnogram.

The remainder of the paper reviews related work (\Cref{sec:related}), fixes
notation, the CP background and the method (\Cref{sec:method}), describes the
experimental setup (\Cref{sec:setup}), reports results (\Cref{sec:results})
and discusses implications, limitations and conclusions
(\Cref{sec:discussion}).
